\section{Introduction}
Educational institutions increasingly rely on browser-based learning platforms to deliver course content to students who use a wide variety of personal devices. While such platforms have matured for text, video, and 2D-interactive content, the integration of interactive 3D visualization remains an unresolved engineering problem. Educational 3D content has well-documented pedagogical value~\cite{Li2021,Yang2024} and established technology pipelines~\cite{Rodrguez2021,Hensen2023}, yet content authors building 3D learning scenarios for institutional platforms face a recurring set of practical obstacles that current rendering solutions do not adequately address.

The educational deployment context introduces specific constraints absent from general-purpose 3D rendering. Students access learning content from heterogeneous personal hardware -- older laptops with integrated graphics, mid-range gaming laptops, school computer-lab workstations of varying age -- over which the institution has no control. Installation of native software is typically not permitted on managed school devices, and even where permitted, students may lack administrative privileges or technical confidence. The content must therefore run inside a standard browser tab, with no plug-ins, on whatever device the student happens to use. Course content teams, in turn, are typically composed of subject-matter experts and software engineers without specialized graphics-programming experience; the cost of building each new 3D learning scenario is dominated by engineering time rather than artistic content.

The two dominant solution paths each fail at one of these constraints. Native game engines such as Unity provide a productive authoring experience and a familiar component-based programming model -- but their browser deployment relies on exporting a full native runtime to WebAssembly. The resulting bundles routinely exceed several megabytes of compressed JavaScript and require multi-second loading times before the first interactive frame, which conflicts with the bandwidth and patience profile of a typical classroom session. Conversely, low-level graphics libraries such as Three.js or Babylon.js produce small and fast-loading applications, but at the cost of pushing the full complexity of rendering architecture~\cite{Evans2014} -- scene graphs, materials, lighting, asset pipelines, frame loops -- onto the application developer. For an educational platform team without dedicated graphics engineers, building and maintaining a scenario directly on Three.js represents a substantial recurring cost.

The educational-platform community therefore lacks a middle-tier solution: a 3D engine that provides the productive, component-based authoring experience of a mature game engine, while remaining lightweight enough to load in a few hundred milliseconds inside a browser tab on consumer hardware. This is the gap the present work addresses.

The present work introduces and deploys such an engine as the runtime for interactive 3D scenarios within an educational learning platform. The engine is implemented in TypeScript, distributed as an npm library, and provides a Unity-inspired component model -- \texttt{GameObject}, \texttt{Component}, \texttt{Transform}, \texttt{Material}, lifecycle hooks (\texttt{awake}, \texttt{start}, \texttt{update}) -- over a fully encapsulated low-level rendering backend (Three.js). Scenario authors write component-based code against the engine API and never encounter the underlying graphics library directly, isolating their code from backend changes and removing the requirement for graphics-programming expertise.

Beyond the architectural design, this paper investigates a complementary problem: the high-level abstraction must not introduce performance penalties that would compromise practical usability on the consumer hardware typical of educational deployments. Three targeted optimization strategies are introduced, addressing distinct dimensions of runtime performance: GPU-compressed texture delivery via the KTX2/Basis Universal format to reduce GPU memory consumption; shader pre-compilation to eliminate first-frame rendering stalls; and dirty-flag transform propagation to minimize per-frame CPU overhead. Cross-device measurements on a mid-range laptop and a desktop workstation show that these optimizations deliver their largest benefits precisely on the lower-tier hardware that students are most likely to use, while imposing no measurable cost on more capable systems.

The contributions of this paper are as follows:
\begin{enumerate}
    \item An engine architecture for educational 3D scenarios that hides a low-level WebGL library behind a Unity-compatible high-level API, including a component model, a scene graph, a material system, a ZIP-based scenario packaging format, and integration with a web-based educational platform.
    \item Three implemented and evaluated optimization strategies targeting the runtime performance dimensions most relevant to consumer educational hardware: asset memory consumption, rendering initialization latency, and per-frame computational overhead.
    \item A cross-device quantitative evaluation demonstrating that the optimization stack delivers a 60\% reduction in maximum frame time and a 97\% reduction in frame-time variability on a mid-range laptop, while a high-end desktop achieves stable 60 FPS regardless of optimization configuration -- the optimizations target the low-end of the hardware spectrum that institutional deployments cannot avoid.
\end{enumerate}

The remainder of this paper is organized as follows. Section~2 reviews related work on browser-based 3D engines, educational 3D platforms, and GPU optimization techniques. Section~3 describes the engine architecture. Section~4 details the three optimization strategies. Section~5 presents the experimental evaluation. Section~6 discusses the implications for educational deployment, and Section~7 concludes with directions for future work.